\section{Human-in-the-Loop}

The concept of \gls{hitl} encapsulates a paradigm where human expertise, judgment, and oversight are integrated into machine learning workflows to enhance model performance, interpretability, and ethical alignment. In the context of LLMs, HITL systems leverage human feedback at various stages: training, evaluation, and deployment, to iteratively refine outputs and ensure they reflect human values, contextual appropriateness, and factual accuracy \cite{chaudhari2024rlhfdecipheredcriticalanalysis}.  
Unlike traditional fully automated training pipelines, HITL approaches recognize that human cognition and ethical reasoning cannot be entirely captured by statistical optimization. Thus, feedback loops serve as a bridge between algorithmic intelligence and human judgment, allowing models to evolve dynamically in response to real-world user interactions \cite{ HelpfulHarmlessHonest2025}.
  

\subsection{Iterative Refinement}

Iterative refinement in HITL workflows with LLMs agent orchestration refers to a repeated generate–evaluate–update loop in which one or more LLM agents produce candidate outputs, a human provides feedback, and the system updates subsequent outputs or orchestration policies accordingly. This process enables continuous improvement through successive correction and alignment cycles, leveraging human expertise to guide LLM behavior toward desired outcomes \cite{wang2022hitl_survey}. 

In the context of agentic systems, iterative refinement can be implemented at different levels: (i) the generation level, where humans correct or critique model outputs; (ii) the policy level, where feedback informs task decomposition or tool selection; and (iii) the orchestration level, where human input adjusts coordination among multiple agents. Recent work such as \textit{Self-Refine} demonstrates the effectiveness of self-feedback and iterative improvement without additional supervised training, allowing models to iteratively critique and improve their own outputs \cite{madaan2023selfrefine}. Similarly, systems like IMPROVE show how iterative, component-wise pipeline refinement using LLM agents can yield stable performance improvements in complex multi-agent pipelines \cite{xue2025improve}.

Iterative refinement may involve several workflow patterns: (1) \textit{Human-as-gate}, where humans approve or reject candidate outputs; (2) \textit{Human-as-critic}, where humans provide structured feedback; (3) \textit{Human-as-instructor}, where humans directly reformulate the task or constraints; and (4) \textit{Automated self-refinement with human validation}, where LLM agents refine their outputs autonomously before human review. The combination of automated self-correction and human oversight has been shown to improve factuality, alignment, and interpretability in complex agent systems \cite{llmhas2025survey, zhang2025_clinical_iterative}.

\subsection{Benefits and Challenges}

The integration of iterative refinement within HITL orchestration presents a number of tangible benefits across dimensions of quality, safety, interpretability, and efficiency. One of the primary advantages lies in its capacity to improve alignment and factual accuracy. By incorporating human evaluation and corrective feedback into iterative loops, systems can mitigate common LLM limitations such as hallucination and task misalignment \cite{wang2022hitl_survey}. The presence of human oversight ensures that model outputs are not only syntactically correct but also contextually appropriate and aligned with human values or domain-specific standards. This quality assurance mechanism is particularly valuable in high-stakes environments, such as healthcare, finance, or law, where factual precision and interpretability are critical \cite{zhang2025_clinical_iterative}.

A further benefit of iterative refinement is its contribution to the overall safety and reliability of agentic systems. When human experts are embedded within the loop, they can intercept unsafe, biased, or inappropriate outputs before they propagate through the system or reach end-users. This “safety valve” effect transforms humans from passive evaluators into active participants in system governance, creating a dynamic safeguard that complements automated safety mechanisms. Moreover, iterative human feedback supports transparency and traceability, as each cycle of revision can be logged and attributed to either human or model actions, which is essential for auditing and compliance purposes. From an operational standpoint, hybrid refinement also enables efficient division of labor between human and artificial agents: routine, low-risk refinements can be automated through self-refinement strategies, while human attention is reserved for complex or ambiguous cases \cite{madaan2023selfrefine, xue2025improve}. In this way, the orchestration of human and LLM agents facilitates both scalability and adaptability, achieving better outcomes than purely manual or fully autonomous systems.

Despite these advantages, HITL orchestration also introduces significant challenges that must be addressed to ensure practical viability and ethical soundness. The most immediate obstacle concerns the cost and latency of human involvement. Continuous human review is resource-intensive, potentially reducing system throughput and limiting scalability in real-world applications \cite{wang2022hitl_survey}. Designing adaptive orchestration strategies that determine when and where human input is most valuable thus becomes a central research question. Another challenge arises from the scalability and coordination of multi-agent systems: as the number of agents increases, orchestrating feedback, tracking dependencies, and maintaining coherence across iterative refinements become increasingly complex \cite{llmhas2025survey}. Ensuring consistent integration of human feedback across multiple agents or pipelines is technically demanding and often requires sophisticated scheduling and communication protocols.

Furthermore, the quality of human feedback itself is not guaranteed. Human annotators or reviewers may introduce bias, inconsistency, or noise into the iterative loop, which can destabilize the learning process if not properly managed \cite{wang2022hitl_survey}. To mitigate these effects, research has explored structured feedback formats, inter-annotator agreement measures, and aggregation techniques. Another critical issue concerns evaluation: the involvement of humans blurs the boundary between model autonomy and external assistance, complicating efforts to measure true model capability or the causal effect of human intervention. Metrics that account for human effort, latency, and risk are still under development and remain an open area of study. Finally, privacy and governance concerns pose additional constraints. When human reviewers handle sensitive data, as in medical or financial domains, organizations must implement strict data protection, access control, and minimization policies \cite{zhang2025_clinical_iterative}. These requirements can limit the applicability of HITL workflows in regulated contexts.

